% Physical unit 025; accepted Session 03 after fresh 400-dpi authority replay.
\GordanUnitStart{025}{printed page 20}
\GordanTopFolio{20}
\begingroup\small
\[
\begin{array}{ll}
\text{aut } l<l_1, & \psi>\omega,\\
\text{aut } l>l_1, & \psi>\omega,\\
\text{aut } l<l_1, & \psi<\omega,\\
\text{aut } l>l_1, & \psi<\omega.
\end{array}
\]
Aequationes linearum primi et quarti generis formam praestant hanc:
\[
 \psi-\omega=\int_{\varphi_1}^{\varphi}\cos l_0\,d\varphi\,{\sqrt{1-e^2+e^2\sin^2l_0\cos^2\varphi}\over 1-\sin^2l_0\cos^2\varphi}
\]
secundi et tertii hanc:
\[
 \omega-\psi=\int_{\varphi_1}^{\varphi}\cos l_0\,d\varphi\,{\sqrt{1-e^2+e^2\sin^2l_0\cos^2\varphi}\over 1-\sin^2l_0\cos^2\varphi}.
\]

\begin{center}\S\ 3.\end{center}

Omnium autem linearum geodeticarum ratio est haec, ut unaquaeque omnia puncta zonae, sphaeroidis circulis $l_0$ et $-l_0$ inclusae contineat.

\begin{center}ARGUMENTATIO.\end{center}

Demonstrandum est, quamvis lineam geodeticam quemlibet circulum zonae deinceps omnibus punctis secare; longitudines igitur eorum, omnes deinceps numeros, intra $0$ et $2\Pi$ sitos, induere.

Quae sectionis longitudines petuntur ex aequatione:
\[
 \psi-\omega=\pm\int_{\varphi_1}^{\varphi}\cos l_0\,d\varphi\,{\sqrt{1-e^2+e^2\sin^2l_0\cos^2\varphi}\over 1-\sin^2l_0\cos^2\varphi}
\]
in qua $\varphi_1$ constans est, valores autem anguli $\varphi$ deducuntur ex aequatione:
\[
 \sin l=\sin l_0\cos\varphi.
\]
Posito in primo sectione:
\[
 \varphi=\varphi_2\qquad \psi=\psi_2,
\]
aliis in sectionibus deinceps fiet: angulus
\[
 \varphi=\varphi_2+2\pi\qquad \varphi_2+4\pi
\]
\[
 \psi=\omega+\int_{\varphi_1}^{\varphi_2+2\pi}\cos l_0\,d\varphi\,{\sqrt{1-e^2+e^2\sin^2l_0\cos^2\varphi}\over 1-\sin^2l_0\cos^2\varphi}
\]
\[
 \omega+\int_{\varphi_1}^{\varphi_2+4\pi}\cos l_0\,d\varphi\,{\sqrt{1-e^2+e^2\sin^2l_0\cos^2\varphi}\over 1-\sin^2l_0\cos^2\varphi}
\]
vel $\omega+4\mathcal A(l_0)$ \qquad $\omega+8\mathcal A(l_0)$
\endgroup
